熵回答的是一个分布自身有多不确定。可实际工作里要比的往往是两个分布：数据按 \(P\) 产生而模型相信 \(Q\)，训练集与线上流量是否已经漂开，两种生成机制值不值得区分。逐点相减、比较参数，这些做法都能给出一个数，但那个数说不清哪一种错误更贵。

这一章从一个具体冲突进入：数据来自 \(P\)，你却按 \(Q\) 分配码长，平均每个符号多付多少比特？这个问题自然长出交叉熵和 KL 散度，也让 KL 的方向性和支撑条件从需要背诵的性质变成了可以推出来的后果。

\begin{center}
\begin{tikzpicture}[x=1cm,y=1cm,font=\small]
\begin{scope}
  \draw[black!55] (0,0.9) circle (0.34) node {\(P\)};
  \draw[black!55] (2.8,0.9) circle (0.34) node {\(Q\)};
  \draw[->,black!65] (0.36,1.12) -- (2.42,1.12);
  \draw[->,black!65] (2.44,0.68) -- (0.38,0.68);
  \node[above,font=\footnotesize,black!75] at (1.4,1.16) {\(\KL(P\|Q)\)};
  \node[below,font=\footnotesize,black!75] at (1.4,0.64) {\(\KL(Q\|P)\)};
  \node[font=\footnotesize,align=center,black!70] at (1.4,-0.35) {两条箭头读数不同，\\有时一条有限一条无穷};
\end{scope}
\begin{scope}[shift={(6.0,0)}]
  \draw[black!55] (0,0.9) circle (0.34) node {\(P\)};
  \draw[black!55] (2.8,0.9) circle (0.34) node {\(Q\)};
  \fill[black!8] (1.4,2.0) circle (0.34);
  \draw[black!55] (1.4,2.0) circle (0.34) node {\(M\)};
  \draw[->,black!65] (0.24,1.18) -- (1.12,1.82);
  \draw[->,black!65] (2.56,1.18) -- (1.68,1.82);
  \node[font=\footnotesize,align=center,black!70] at (1.4,-0.35) {改与共同混合比较，\\对称、有界、不发散};
\end{scope}
  \draw[black!25] (4.2,-0.75) -- (4.2,2.5);
\end{tikzpicture}

\emph{左边是 KL 的处境：两个分布互相充当参照，方向一换读数就变，参照分布判零的地方还会直接发散。右边是 Jensen--Shannon 的补救：先造一个双方都覆盖得住的混合分布 \(M\)，再各自与它比较。}
\end{center}

三个判断值得先记下来。其一，\(\KL(P\|Q)\) 的含义是用错模型的代价，而不是分布之间的距离；账单按 \(P\) 的权重结算，所以它天然不对称。其二，两个方向对应完全不同的拟合行为——一个摊开去覆盖，一个缩进某个峰里——这条直接解释了变分推断为什么低估方差、生成模型为什么或者模糊或者只会画同一张脸。其三，交叉熵等于熵加 KL，因此分类训练降低损失就是在降低 KL，损失曲线压不下去的那条渐近线是数据自身的熵。

\subsection{本章篇目}

\begin{itemize}
\tightlist
\item \hyperref[sec:07-Information-Theory/03-Divergences/01]{``错误模型的代价：从交叉熵到 KL 散度''}：把``模型错了''换算成每个符号多付的比特，并指出交叉熵损失与 KL 只差一个常数；
\item \hyperref[sec:07-Information-Theory/03-Divergences/02]{``Gibbs 不等式与 KL 的方向性''}：用 Jensen 一行锁死账单非负，然后把两个方向的行为差别画透；
\item \hyperref[sec:07-Information-Theory/03-Divergences/03]{``Jensen--Shannon 散度与尺度选择''}：支撑不重合时 KL 失效，混合参照怎样救场，以及什么时候该换成 Wasserstein。
\end{itemize}

第一次读，把 KL 读成账单就够用了：每次见到 \(\KL(P\|Q)\)，先说清谁产生数据、谁充当参照，再检查参照有没有在真会发生的地方判零。想弄清原理，再去追 Gibbs 不等式的等号条件、方向性对模型族的挑选，以及混合分布如何换来对称与有界。

前一章是\hyperref[ch:055]{``联合信息：一个变量能告诉我们另一个变量多少''}，那里的互信息在本章会以 \(\KL(P_{XY}\|P_XP_Y)\) 的身份重新出现一次——两条看似不同的线索，其实是同一个不等式的两种写法。
